\documentclass[11pt]{article}
\usepackage[a4paper,margin=25mm]{geometry}
\usepackage[T1]{fontenc}
\usepackage{amsmath,amssymb,amsthm,mathtools,mathrsfs}
\usepackage[hidelinks]{hyperref}
\setlength{\emergencystretch}{2em}
\newtheorem{theorem}{Theorem}[section]
\newtheorem{proposition}[theorem]{Proposition}
\newtheorem{lemma}[theorem]{Lemma}
\newtheorem{corollary}[theorem]{Corollary}
\newcommand{\C}{\mathbb C}\newcommand{\R}{\mathbb R}
\newcommand{\PP}{\mathcal P}\newcommand{\HH}{\mathcal H}
\newcommand{\B}{\mathscr B}\newcommand{\M}{\mathcal M}
\newcommand{\norm}[1]{\left\lVert#1\right\rVert}
\newcommand{\ac}[2]{\{#1,#2\}}
\title{The native arithmetic prolate action, its Lie defects,\\
and the exact polynomial-quotient receiver}
\author{}\date{20 September 2026}
\begin{document}\maketitle

\section{Unchanged arithmetic measure, mass, and physical source}

Retain the original complete finite divisor
\[
 h(s)=\prod_{\rho\in Z}(s-\rho)^{m_\rho},\qquad
 g(s)=2\xi(s)=s(s-1)\pi^{-s/2}\Gamma(s/2)\zeta(s).
\]
Here $Z$ is the original nonempty actual set of nontrivial zeros,
closed under conjugation and $s\mapsto1-s$, and every $m_\rho$
is its complete order. This stipulates the original actual packet;
it does not assert the existence of an off-line packet. For its
unchanged tensor degree $k\geq1$ put
\begin{equation}\label{eq:measure}
 \begin{gathered}
 w_h(u)=\frac{|(g/h)(1/2+iu)|^2}{2\pi},\qquad
 m(u)=m_{h,k}(u)=w_h^{*k}(u),\\
 \mu_h=\int_\R w_h(u)\,du,\qquad
 M=\int_\R m(u)\,du=\mu_h^k,\qquad s_0=k/2,\qquad S=s_0+iu.
 \end{gathered}
\end{equation}
The Hilbert space in this note is the actual
$\HH_m=L^2(\R,m(u)\,du)$, with
$\langle f,g\rangle_m=\int\overline f g\,m\,du$.
No probability density is substituted for $m$.

The density is even and positive almost everywhere. Indeed
$g(\bar s)=\overline{g(s)}$, and the same is true of $h$,
so $w_h$ is even. The zeros of the nonzero entire function
$g/h$ are discrete, hence $w_h>0$ almost everywhere.
Convolution preserves evenness and positivity almost everywhere,
and Tonelli gives the exact mass in \eqref{eq:measure}.
For every $0<t<\pi/2$,
\begin{equation}\label{eq:expmoment}
 \int_\R e^{t|u|}m(u)\,du<\infty.
\end{equation}
Here is a direct sufficient bound. On $\Re s=1/2$ the integral
representation of the alternating series gives
$|\zeta(s)|\leq2|s|/(\sqrt2-1)$: use
$\eta(s)=s\int_1^\infty A(x)x^{-s-1}\,dx$, $|A(x)|\leq1$,
and $\eta(s)=(1-2^{1-s})\zeta(s)$.
The vertical-line Gamma estimate gives a polynomial times
$e^{-\pi|u|/4}$ for $g(1/2+iu)$. Outside a fixed compact interval
$|h(1/2+iu)|$ is bounded below by a positive constant times
$(1+|u|)^{\deg h}$. On that compact interval the quotient
$g/h$ has its removable values and is bounded. Squaring proves
exponential integrability of $w_h$ for $t<\pi/2$.
Since $e^{t|\sum u_j|}\leq\prod_j e^{t|u_j|}$, Tonelli proves
\eqref{eq:expmoment} for its actual convolution.

The original physical source has $D=-x\partial_x$,
$\Theta\phi(x)=2\sum_{n\geq1}\phi(nx)$, and
$\phi_*(x)=(4\pi^2x^4-6\pi x^2)e^{-\pi x^2}$.
Let $F_h\in\B$ be its original divided function,
$\M F_h=g/h$, where $\M F(s)=\int_0^\infty F(x)x^{s-1}\,dx$
and $\B$ has all seminorms
$\sup_{x>0}x^b|D^jF(x)|$, $b\in\mathbb Z$, $j\geq0$.
For completeness, the Fourier transform with kernel
$e^{-2\pi ixy}$ fixes $\phi_*$: differentiating the Gaussian
Fourier identity twice and four times gives
$\widehat{\phi_*}=\phi_*$. Also $\phi_*(0)=\int_\R\phi_*=0$.
Poisson summation therefore gives
$\Theta\phi_*(x)=x^{-1}\Theta\phi_*(1/x)$.
The Gaussian bound at infinity and this identity at zero,
after every $D$ derivative, prove $\Theta\phi_*\in\B$.
The elementary Gaussian Mellin integral and the Gamma
recurrence give
$\M(2\phi_*)(s)=s(s-1)\pi^{-s/2}\Gamma(s/2)$.
Termwise integration on $\Re s>1$ and analytic continuation
then give $\M(\Theta\phi_*)=g$ with every constant retained.
This construction can be made explicitly: starting from
$\Theta\phi_*$, divide once at each root occurrence using
$S_\rho F(x)=x^{-\rho}\int_x^\infty F(y)y^{\rho-1}\,dy$.
The vanishing Mellin moment rewrites this as the negative
integral from zero to $x$, proving all rapid endpoint bounds;
$(D-\rho)S_\rho F=F$ and integration by parts give the Mellin
division. Thus the complete multiplicity produces this
unique $F_h$, with no change of physical coordinate.

For $P\in\C[S]$ define the unchanged tensor source
\begin{equation}\label{eq:physical}
 \mathcal A P=P(D_1+\cdots+D_k)F_h^{\otimes k}.
\end{equation}
Its multiple Mellin transform is
$P(s_1+\cdots+s_k)\prod_j(g/h)(s_j)$.
Mellin Parseval on each original $dx_j$ axis, followed by
Tonelli and convolution, therefore proves
\begin{equation}\label{eq:isometry}
 \norm{\mathcal A P}_{L^2(dx_1\cdots dx_k)}^2
  =\int_\R|P(s_0+iu)|^2m(u)\,du.
\end{equation}
All polynomial pairings follow by polarization.
In particular the constant-source squared norm is $M$, not one.
The finite polynomial maps also retain the original strong
source topology, independently of this Hilbert identity.
On the original $k$-variable rapid space write
\[
 p_{\boldsymbol b,\boldsymbol d}(F)
 =\sup_{\boldsymbol x>0}
       \left|\prod_{\ell=1}^k x_\ell^{b_\ell}
                    D_\ell^{d_\ell}F(\boldsymbol x)\right|.
\]
For $P(S)=\sum_{j=0}^L p_jS^j$ the multinomial theorem gives
the explicit finite strong bound
\begin{equation}\label{eq:strongfinite}
 p_{\boldsymbol b,\boldsymbol d}(\mathcal A P)
 \leq\sum_{j=0}^L|p_j|
       \sum_{\substack{\boldsymbol\alpha\geq0\\|\boldsymbol\alpha|=j}}
       \frac{j!}{\prod_\ell\alpha_\ell!}
       \prod_{\ell=1}^k
                  p_{b_\ell,d_\ell+\alpha_\ell}(F_h).
\end{equation}
Thus every finite coefficient map used below is a continuous
map into the unchanged strong physical source.
The Hilbert completion constructed below instead has
codomain the physical Hilbert space; it is not asserted
to take every Hilbert vector into the rapid source.

\section{Native orthogonal polynomials and exact coefficient maps}

Let $Q_n(u)$ be the unique real monic orthogonal polynomials for
$m(u)\,du$, and retain their full norms
\begin{equation}\label{eq:norms}
 \omega_n=\int_\R Q_n(u)^2m(u)\,du>0,\qquad
 \omega_0=M,\qquad
 \phi_n=Q_n/\sqrt{\omega_n},\qquad
 a_n=\sqrt{\omega_{n+1}/\omega_n},\quad a_{-1}=0.
\end{equation}
Positive density makes every finite moment Gram positive
definite, proving existence and uniqueness by Gram--Schmidt.
Evenness and uniqueness give $Q_n(-u)=(-1)^nQ_n(u)$.
The exact recurrence is
\begin{equation}\label{eq:recurrence}
 uQ_n=Q_{n+1}+\frac{\omega_n}{\omega_{n-1}}Q_{n-1}
 \quad(n\geq1),\qquad uQ_0=Q_1;
 \qquad
 u\phi_n=a_{n-1}\phi_{n-1}+a_n\phi_{n+1}.
\end{equation}
To verify the recurrence, pair $uQ_n$ with $Q_j$ for $j<n-1$
and move $u$ to $Q_j$, whose degree is at most $n-1$.
The diagonal pairing vanishes by parity. The leading
coefficient fixes $Q_{n+1}$, and pairing with $Q_{n-1}$
gives $\omega_n/\omega_{n-1}$.

The $\phi_n$ are complete in $\HH_m$. If $f$ is orthogonal
to every polynomial, Cauchy--Schwarz and \eqref{eq:expmoment}
make $\int\overline f(u)e^{zu}m(u)\,du$ holomorphic on a
strip containing the imaginary axis. Every derivative at
zero vanishes, so the function vanishes on that strip.
Fourier uniqueness for the finite measure $\overline f m\,du$
gives $f=0$. Hence the map
\[
 U:\ell^2(\mathbb N_0)\longrightarrow\HH_m,\qquad Ue_n=\phi_n
\]
is unitary. The multiplication operator by $u$ on its natural
domain corresponds to the Jacobi operator $J$,
$Je_n=a_{n-1}e_{n-1}+a_ne_{n+1}$ on the finite-support core.
Indeed the same exponential-moment density argument applied
to $(1+u^2)m(u)\,du$ proves density of polynomials in the
graph norm of multiplication by $u$, so this is a core.
The number operator $\mathsf N e_n=ne_n$ is self-adjoint
on $\{\sum x_ne_n:\sum n^2|x_n|^2<\infty\}$.
All products and brackets below are asserted on the
invariant polynomial/finite-support core; no self-adjoint
extension of a prolate expression is selected by them.

The original monic $S$-polynomials, including their phases, are
\begin{equation}\label{eq:physical-polys}
 p_n(S)=i^nQ_n((S-s_0)/i),\qquad
 p_n(s_0+iu)=i^nQ_n(u),\qquad \norm{p_n}_m^2=\omega_n.
\end{equation}
The constant vector $1$ is $\sqrt M\,\phi_0$.
The corresponding unit cyclic vector, when the convention of
a unit cyclic pair is wanted, is explicitly $1/\sqrt M$
in this same measure; the physical vector $1$ and its mass
are never replaced.

Here are full finite coordinate maps. Let $B_L$ have as its
$n$th column the $u$-power coefficients of $\phi_n$, $0\leq n\leq L$.
Let
\begin{equation}\label{eq:coordinates}
 (T_L)_{rj}=\binom jr s_0^{j-r}i^r\quad(r\leq j),\qquad
 (T_L)_{rj}=0\quad(r>j),\qquad
 D_L=B_L^{-1}T_L.
\end{equation}
Since the diagonal of $B_L$ is $1/\sqrt{\omega_n}$,
it is invertible, and
$P(s_0+iu)=\sum_{n=0}^L(D_L\operatorname{coeff}P)_n\phi_n(u)$.
Thus the complete original source Gram in increasing $S$ powers is
\begin{equation}\label{eq:sourcegram}
 H_{S,L}=D_L^*D_L.
\end{equation}
The native grading in those exact coordinates is
$D_L^{-1}\operatorname{diag}(0,\ldots,L)D_L$.
The rectangular map of multiplication by $S$ from degree
$\leq L$ to degree $\leq L+1$ is the corresponding rectangular
matrix $s_0I+iJ$, with the outgoing $a_L$ edge included.
These equalities prove all physical phases and mass factors,
not just equivalence of the two polynomial spans.

Equation \eqref{eq:isometry} extends the map
$f(u)\mapsto\mathcal A(f((S-s_0)/i))$ from polynomials to
a unitary map from $\HH_m$ onto the closed cyclic Hilbert
subspace generated by the original tensor vector.
It does not identify that subspace with the whole tensor
Hilbert space. On its polynomial domain it transports
$-iJ$ to $s_0-(D_1+\cdots+D_k)$ exactly.

\section{The prolate polynomial action and the phase dictionary}

Connes--Consani--Moscovici \cite{CCM}, Definition
\texttt{prolop}, attach the formal expression $-J^2+\eta^2\mathsf N$
to an even cyclic pair. Their actual archimedean prolate
formula in Theorem \texttt{mainstart}(v) is instead
$-J^2+2\pi\lambda^2(4\mathsf N+1)-1/4$.
Retaining all those constants, define for real $\lambda$
its native arithmetic receiver
\begin{equation}\label{eq:W}
 W_\lambda=-J^2+8\pi\lambda^2(\mathsf N+1/4)-1/4
 \quad\text{on the finite-support core}.
\end{equation}
This is a densely defined symmetric operator.
It is a polynomial-domain receiver, not an assertion about
a self-adjoint prolate extension or its spectrum.
Its entries are, with missing negative indices zero,
\begin{equation}\label{eq:Wentries}
 \begin{split}
 W_\lambda e_n={}&-a_{n-1}a_{n-2}e_{n-2}
  +w_ne_n-a_na_{n+1}e_{n+2},\\
 w_n={}&-a_n^2-a_{n-1}^2+8\pi\lambda^2(n+1/4)-1/4.
 \end{split}
\end{equation}
These follow by applying \eqref{eq:recurrence} twice.
The parity involution $\Gamma e_n=(-1)^ne_n$ commutes with
$W_\lambda$ and $\mathsf N$ and anticommutes with $J$.
Consequently on the parity-$\varepsilon$ sector
($n=2r+\varepsilon$, $\varepsilon=0,1$), the adjacent off-diagonal
is $-a_{2r+\varepsilon}a_{2r+\varepsilon+1}$ and the
diagonal is $w_{2r+\varepsilon}$.

Let $\mathcal N$ be the polynomial grading
$\mathcal Np_n=np_n$ in the original $S$ coordinate.
The physical polynomial-domain action of \eqref{eq:W} is
\begin{equation}\label{eq:Wphysical}
 (\mathcal W_\lambda P)(S)
 =(S-s_0)^2P(S)+8\pi\lambda^2(\mathcal N+1/4)P(S)-P(S)/4.
\end{equation}
It maps degree $\leq L$ into degree $\leq L+2$.
Its actual source vector is $\mathcal A\mathcal W_\lambda P$.
In the orthonormal physical frame $p_n/\sqrt{\omega_n}$,
the two off-diagonals in \eqref{eq:Wentries} have the opposite
sign, namely $+a_na_{n+1}$: the exact diagonal change of
frame is $e_n\mapsto i^ne_n$, and $i^{n}/i^{n+2}=-1$.
This explains the positive off-diagonals in the archimedean
phase convention of \cite{CCM}; no phase is discarded.

\section{All Lie brackets, with the literal generator factors}

Fix a real scalar $c$, distinct from the physical centre $s_0$,
and keep precisely
\begin{equation}\label{eq:generators}
 H=-iJ,\qquad K=\mathsf N+c,\qquad
 B=\frac{[J,\mathsf N]}{2i},\qquad
 E_+=\frac i2(K+B),\qquad E_-=\frac i2(-K+B).
\end{equation}
Define the actual diagonal discrepancy
\begin{equation}\label{eq:delta}
 \mathcal D e_n=(a_n^2-a_{n-1}^2)e_n,\qquad
 \Delta_c=\mathcal D-2K,\qquad
 \delta_n^{(c)}=a_n^2-a_{n-1}^2-2n-2c.
\end{equation}

\begin{theorem}[Complete bracket calculation]\label{thm:brackets}
For every positive native recurrence sequence,
\begin{equation}\label{eq:basicbrackets}
 [H,K]=2B,\qquad [H,B]=2K+\Delta_c,\qquad [K,B]=-H/2.
\end{equation}
Consequently the literal three generators satisfy
\begin{equation}\label{eq:literalbrackets}
 \boxed{
 [H,E_+]=2E_++\tfrac i2\Delta_c,\quad
 [H,E_-]=-2E_-+\tfrac i2\Delta_c,\quad
 [E_+,E_-]=H/4.}
\end{equation}
The remaining mixed brackets are
$[K,E_\pm]=-iH/4$, $[B,E_+]=iH/4$,
$[B,E_-]=-iH/4$, with reverse brackets their negatives.
In particular the literal standard-$\mathfrak{sl}_2$
third-bracket defect is $-3H/4$, even if $\Delta_c=0$.

The first two weight relations hold without defect if and only if
\begin{equation}\label{eq:allc}
 a_n^2=(n+1)(n+2c)\quad(n\geq0).
\end{equation}
Positivity is equivalent to $c>0$ within this family.
For every such $c$, the explicitly different generators
\begin{equation}\label{eq:morphism}
 F_+=2E_+,\qquad F_-=2E_-,\qquad K_{\rm std}=2K
\end{equation}
satisfy
$[H,F_\pm]=\pm2F_\pm$, $[F_+,F_-]=H$,
$i(F_--F_+)=K_{\rm std}$.
These are an exact change of Lie-algebra basis, not an
unannounced change in \eqref{eq:generators}.
\end{theorem}
\begin{proof}
Put $Re_n=a_ne_{n+1}$ and $Le_n=a_{n-1}e_{n-1}$.
Then $J=L+R$, $[J,\mathsf N]=L-R$,
$B=(L-R)/(2i)$, and $[L,R]=\mathcal D$.
Also $[\mathsf N,L]=-L$ and $[\mathsf N,R]=R$.
Direct substitution gives $[J,K]=2iB$,
$[J,B]=i\mathcal D$, $[K,B]=iJ/2$, proving
\eqref{eq:basicbrackets}. For example
$[H,E_+]=(i/2)(2B+2K+\Delta_c)
 =2E_++i\Delta_c/2$; the negative case is the same
substitution with the sign on $K$ reversed.
Finally
\[
 [E_+,E_-]=-\tfrac14[K+B,-K+B]
          =-\tfrac12[K,B]=H/4.
\]
This verifies the third bracket independently of any
recurrence hypothesis. The mixed brackets follow by the
same two-term expansions.

The vanishing of the two weight defects is $\Delta_c=0$.
Sum $a_n^2-a_{n-1}^2=2(n+c)$ from $0$ to $n$, with
$a_{-1}=0$, to obtain \eqref{eq:allc}. Conversely taking
the difference of its two consecutive values gives
the required diagonal equality. Its first value is
$a_0^2=2c$, and all values are positive exactly for $c>0$.
Multiplying both $E$'s by two proves every identity
in \eqref{eq:morphism}.
\end{proof}

This computation supplies a precise correction at the scope
of \cite{CCM}. It applies to equations \texttt{h} and
\texttt{eplusminusricu}, and to Theorem \texttt{dzero}(i).
The asserted uniqueness of $c=1/4$ does not follow from the
written relations. With the literal generators there is
also the third-bracket factor in \eqref{eq:literalbrackets}.
The archimedean even representation has an additional
specified compact spectrum $2n+1/2=2(n+1/4)$, which fixes
$c=1/4$ after the exact map \eqref{eq:morphism}.
That extra representation datum is not a theorem about
the native arithmetic recurrence.

\begin{proposition}[Quadratic operator and exact Casimir receiver]
\label{prop:casimir}
The standard-basis quadratic expression, retaining the
literal $E$'s, is
\begin{equation}\label{eq:Cas}
 \begin{split}
 \mathcal Q_c
 &=H^2+8(E_+E_-+E_-E_+)
   =4K^2-J^2-4B^2,\\
 \mathcal Q_ce_n
 &=\bigl(4(n+c)^2-2(a_n^2+a_{n-1}^2)\bigr)e_n.
 \end{split}
\end{equation}
For arbitrary coefficients its commutator defects are
\begin{equation}\label{eq:Casdef}
 [\mathcal Q_c,H]=4\ac B{\Delta_c},\quad
 [\mathcal Q_c,K]=0,\quad
 [\mathcal Q_c,B]=\ac H{\Delta_c},\quad
 [\mathcal Q_c,E_\pm]=\tfrac i2\ac H{\Delta_c}.
\end{equation}
Here $\ac AB=AB+BA$. Under \eqref{eq:allc},
$\mathcal Q_c=4c(c-1)I$, which is $-3I/4$ at $c=1/4$.
The unmodified expression
$C=H^2+2(E_+E_-+E_-E_+)$ instead obeys
$C=(\mathcal Q_c+3H^2)/4$; even under
\eqref{eq:allc} it is not central.
\end{proposition}
\begin{proof}
Expand $E_+E_-+E_-E_+=(K^2-B^2)/2$.
Then $J^2+4B^2=2(LR+RL)$, which proves the diagonal
formula. Expanding commutators of squares using
\eqref{eq:basicbrackets} gives, for example,
\[
 [\mathcal Q_c,H]
  =-8\ac KB+4\ac B{2K+\Delta_c}
  =4\ac B{\Delta_c}.
\]
The $K$ terms cancel; the $B$ calculation is
$\ac H{2K+\Delta_c}-2\ac KH=\ac H{\Delta_c}$.
The $E$ equations follow by linearity.
For \eqref{eq:allc},
$a_n^2+a_{n-1}^2=2n^2+4cn+2c$,
so the diagonal is $4c(c-1)$.
The equation for $C$ follows by direct subtraction.
Its noncentrality is explicit:
$[C,K]=(3/2)\ac HB=(3/2)(R^2-L^2)$,
whose value on $e_0$ is $(3/2)a_0a_1e_2\ne0$.
\end{proof}

\section{All masses in the zero-discrepancy family}

The free parameter in \eqref{eq:allc} corresponds to an exact
family of measures, not merely a formal sequence. For $c>0$ let
\begin{equation}\label{eq:gamfamily}
 dm_{c,M}(u)=
 \frac{M\,2^{2c}}{4\pi\Gamma(2c)}
       |\Gamma(c+iu/2)|^2\,du .
\end{equation}
Its mass is $M$ and its moment-generating function is
$M(\cos t)^{-2c}$ for $|t|<\pi/2$.
For a verification with the factors retained, the beta integral
after $x=e^y$ is
\[
 \frac{\Gamma(c+iv)\Gamma(c-iv)}{\Gamma(2c)}
  =\int_\R (2\cosh(y/2))^{-2c}e^{ivy}\,dy .
\]
Fourier inversion, then $u=2v$, gives
$\int|\Gamma(c+iu/2)|^2e^{iux}\,du
 =4\pi\Gamma(2c)2^{-2c}(\cosh x)^{-2c}$.
The Gamma exponential bound permits analytic continuation
to $x=-it$, proving the asserted mass and transform.

Define monic polynomials by
\begin{equation}\label{eq:genfamily}
 \sum_{n\geq0}Q_n^{(c)}(u)\frac{z^n}{n!}
       =(1+z^2)^{-c}e^{u\arctan z}.
\end{equation}
Differentiating in $z$ proves the recurrence
$uQ_n^{(c)}=Q_{n+1}^{(c)}+n(n-1+2c)Q_{n-1}^{(c)}$.
For sufficiently small real $z,w$, multiplying the two
generating functions and integrating against
\eqref{eq:gamfamily} gives
\[
 M(1+z^2)^{-c}(1+w^2)^{-c}
       \bigl(\cos(\arctan z+\arctan w)\bigr)^{-2c}
       =M(1-zw)^{-2c}.
\]
Exponential domination justifies coefficient extraction,
so the full norms are
\begin{equation}\label{eq:gamnorms}
 \int Q_n^{(c)}Q_j^{(c)}\,dm_{c,M}
       =\delta_{nj}M\,n!(2c)_n .
\end{equation}
Thus \eqref{eq:allc} is realized for every $c>0$ with the
same prescribed mass $M$.

For the native arithmetic measure the following classification
is exact for $c>0$:
\begin{equation}\label{eq:exacttest}
 \begin{gathered}
 \Delta_c=0
 \quad\Longleftrightarrow\quad m(u)\,du=dm_{c,M}(u)\\
 \quad\Longleftrightarrow\quad
 \left(\int_\R e^{tu}w_h(u)\,du\right)^k
          =\mu_h^k(\cos t)^{-2c}\quad(|t|<\pi/2).
 \end{gathered}
\end{equation}
Indeed equal recurrences and the same $\omega_0=M$ give
identical moments:
$\int u^r\,dm=M\langle e_0,J^re_0\rangle$ is a finite
weighted-path calculation determined by the coefficients.
Both measures have exponential moments, so equality of
all moments gives equality of their analytic Laplace
transforms and then their Fourier transforms and measures.
The reverse implication follows from \eqref{eq:gamnorms};
the convolution Laplace identity proves the last equivalence.

\begin{theorem}[The native arithmetic Lie defect never vanishes]
\label{thm:nonzero}
For every actual $h$ and $k$ in \eqref{eq:measure} and every
real $c$, the operator $\Delta_c$ is nonzero. In particular,
there is at least one finite index $n$ with
$a_n^2-a_{n-1}^2-2n-2c\ne0$.
\end{theorem}
\begin{proof}
For $c\leq0$, its first entry is
$\delta_0^{(c)}=a_0^2-2c>0$.
Suppose $c>0$ and $\Delta_c=0$.
For real $|t|<\pi/2$ the native Laplace transform is positive.
Taking its unique positive $k$th root in
\eqref{eq:exacttest} gives
\[
 \int_\R e^{tu}w_h(u)\,du
       =\mu_h(\cos t)^{-2c/k}.
\]
Analytic Laplace uniqueness as proved above, with
$d=c/k>0$, forces the exact equality of densities
\begin{equation}\label{eq:impossiblegamma}
 w_h(u)=\frac{\mu_h\,2^{2d}}{4\pi\Gamma(2d)}
                \Gamma(d+iu/2)\Gamma(d-iu/2).
\end{equation}
The equality is first almost everywhere and then everywhere
on the real axis, since both densities are continuous.
But $F=g/h$ is entire, and
\[
 W_h(z)=\frac{F(1/2+iz)F(1/2-iz)}{2\pi}
\]
is an entire function agreeing with $w_h$ on the real axis.
The right side of \eqref{eq:impossiblegamma} is holomorphic
on $|\Im z|<2d$. The identity theorem makes it equal to
$W_h$ throughout that strip. At $z_0=2id$ its meromorphic
continuation has the nonzero residue
\[
 \frac{\mu_h2^{2d}}{4\pi\Gamma(2d)}
       \frac{2}{i}\Gamma(2d)
       =-\frac{i\mu_h2^{2d-1}}{\pi}.
\]
Approaching $z_0$ from inside the strip makes the right
side unbounded, whereas the entire $W_h$ is bounded near
$z_0$. This contradiction proves the theorem.
\end{proof}

Thus the receiving formulas describe an actual nonzero
arithmetic defect, not an unverified defect-free
representation. The theorem does not identify its first
nonzero index; every entry is exactly evaluated by the
original moment determinants below. Nonvanishing alone
gives no growing-degree lower bound and no uniform bound
as $h$, $k$, or the cutoff varies.
In particular, at $c=1/4$ the original arithmetic discrepancy
is exactly
\begin{equation}\label{eq:archdefect}
 \boxed{\quad
 \Delta_{1/4}e_n=
       \left(\frac{\omega_{n+1}}{\omega_n}
        -\frac{\omega_n}{\omega_{n-1}}-2n-\frac12\right)e_n,
 \quad}
\end{equation}
where the second ratio is zero for $n=0$.
Its first entry retains the actual convolution variance:
$\delta_0^{(1/4)}
 =k\bigl(\int u^2w_h(u)\,du/\mu_h\bigr)-1/2$.
No Gamma replacement of that integral is made.

\section{Original Hankel determinants and Toda differences}

Let the unscaled moments and their exponential deformation be
\begin{equation}\label{eq:tau}
 \mu_j(t)=\int_\R u^je^{tu}m(u)\,du,\qquad
 \tau_n(t)=\det(\mu_{r+s}(t))_{0\leq r,s<n},\quad \tau_0(t)=1 .
\end{equation}
These are well-defined and analytic for real $|t|<\pi/2$
and holomorphic locally there. For real $t$ the Grams are
positive definite. At zero $\tau_1(0)=M$, with no mass
cancellation. These are also exactly the original
$S$-coordinate Gram determinants, not replacement volumes:
if
\[
 H_{S,n-1}(t)_{rj}
 =\int_\R\overline{(s_0+iu)^r}(s_0+iu)^j e^{tu}m(u)\,du
 \quad(0\leq r,j<n),
\]
then \eqref{eq:coordinates} gives
$H_{S,n-1}(t)=T_{n-1}^*(\mu_{r+j}(t))T_{n-1}$.
Its triangular coordinate determinant is
$\det T_{n-1}=i^{n(n-1)/2}$, so
$\det H_{S,n-1}(t)=\tau_n(t)$ exactly.
Gram--Schmidt using monic polynomials gives
\begin{equation}\label{eq:hankelratios}
 \omega_n=\frac{\tau_{n+1}(0)}{\tau_n(0)},\qquad
 a_n^2=\frac{\tau_{n+2}(0)\tau_n(0)}{\tau_{n+1}(0)^2}.
\end{equation}
The determinant of the monic triangular change of basis is
one, so the Gram determinant is precisely
$\prod_{j=0}^{n-1}\omega_j$; this proves both formulas.
Therefore every entry of the native Lie defect is the
explicit original Gram expression
\begin{equation}\label{eq:hankeldefect}
 \delta_n^{(c)}
  =\frac{\tau_{n+2}\tau_n}{\tau_{n+1}^2}
    -\frac{\tau_{n+1}\tau_{n-1}}{\tau_n^2}-2n-2c,
 \qquad (n\geq1),
 \quad \delta_0^{(c)}=\frac{\tau_2}{\tau_1^2}-2c,
\end{equation}
where unlabelled $\tau$'s in this display are evaluated at zero.

\begin{theorem}[Exact Toda receiver]\label{thm:toda}
For $n\geq1$ and real $|t|<\pi/2$,
\begin{equation}\label{eq:toda}
 \tau_n\tau_n''-(\tau_n')^2=\tau_{n+1}\tau_{n-1},
 \qquad \partial_t^2\log\tau_n
                    =\frac{\tau_{n+1}\tau_{n-1}}{\tau_n^2}.
\end{equation}
Let $\omega_n(t)=\tau_{n+1}(t)/\tau_n(t)$,
$\beta_n(t)=\partial_t\log\omega_n(t)$ and
$a_n(t)^2=\omega_{n+1}(t)/\omega_n(t)$, with $a_{-1}(t)=0$.
Then
\begin{equation}\label{eq:todarecurrence}
 \begin{gathered}
 \beta_n'(t)=a_n(t)^2-a_{n-1}(t)^2,\qquad
 \partial_t\log a_n(t)^2=\beta_{n+1}(t)-\beta_n(t),\\
 \beta_n(0)=0,\qquad
 \delta_n^{(c)}
   =\left.\partial_t^2\log\frac{\tau_{n+1}(t)}{\tau_n(t)}
                                               \right|_{t=0}-2n-2c .
 \end{gathered}
\end{equation}
\end{theorem}
\begin{proof}
Differentiation under the integral gives
$\mu_j'(t)=\mu_{j+1}(t)$. Differentiate $\tau_n$ by rows:
every differentiated row except the last repeats the next
row, so $\tau_n'$ is the minor with row indices
$0,\ldots,n-2,n$ and column indices $0,\ldots,n-1$.
Differentiate that expression by columns to obtain
$\tau_n''$ with both row and column indices
$0,\ldots,n-2,n$. Apply the Desnanot--Jacobi identity
to the $(n+1)$-square Hankel matrix with indices $0,\ldots,n$,
deleting the last two indices. It gives exactly
$\tau_n\tau_n''-(\tau_n')^2=\tau_{n+1}\tau_{n-1}$.
This identity itself follows by two-row/two-column
Schur elimination on an invertible upper-left block;
that block is positive definite here. Dividing by
$\tau_n^2$ proves the logarithmic formula.
Taking the difference at $n+1$ and $n$ proves the first
identity in \eqref{eq:todarecurrence}; differentiating
the defining norm ratio proves the second.
Evenness makes $\mu_j(-t)=(-1)^j\mu_j(t)$ and thus
$\tau_n(-t)=\tau_n(t)$, since the row and column signs
multiply to $(-1)^{n(n-1)}=1$. Hence $\beta_n(0)=0$.
The last identity is \eqref{eq:delta} with the first
Toda identity substituted.
\end{proof}

The first two discrepancies can therefore be evaluated
without an orthogonalization ambiguity:
\[
 \delta_0^{(c)}=\frac{\mu_2}{\mu_0}-2c,\qquad
 \delta_1^{(c)}=\frac{\mu_4}{\mu_2}
              -2\frac{\mu_2}{\mu_0}-2-2c,
\]
with moments at zero; here
$\tau_2=\mu_0\mu_2$ and
$\tau_3=\mu_2(\mu_0\mu_4-\mu_2^2)$.
The full sequence is retained in \eqref{eq:hankeldefect},
not replaced by these initial entries.

\section{Every finite cutoff, including excursions beyond the edge}

Let $\Pi_L$ project to $e_0,\ldots,e_L$, write
$P_L^\partial=e_Le_L^*$, and put $\alpha_L=a_L^2$.
Compressed single generators carry a subscript $L$.
Because $\mathsf N$ preserves the cutoff,
$B_L=[J_L,\mathsf N_L]/(2i)$ exactly.

\begin{theorem}[Complete finite edge corrections]\label{thm:cutoff}
The finite diagonal discrepancy is
\begin{equation}\label{eq:cutdelta}
 \Delta_{c,L}=\Pi_L\Delta_c\Pi_L-\alpha_LP_L^\partial.
\end{equation}
All identities \eqref{eq:basicbrackets} and
\eqref{eq:literalbrackets} hold with the finite matrices and
this $\Delta_{c,L}$. In particular
\begin{equation}\label{eq:cutbracket}
 [H_L,E_{\pm,L}]
 =\pm2E_{\pm,L}+\tfrac i2\Pi_L\Delta_c\Pi_L
                         -\tfrac i2\alpha_LP_L^\partial,
 \qquad [E_{+,L},E_{-,L}]=H_L/4.
\end{equation}
The exact compressed prolate action is
\begin{equation}\label{eq:cutprolate}
 \boxed{\quad
 W_{\lambda,L}:=\Pi_LW_\lambda\Pi_L
 =-J_L^2-\alpha_LP_L^\partial
        +8\pi\lambda^2(\mathsf N_L+1/4)-1/4.\quad}
\end{equation}
The quadratic expression built from the finite generators
instead satisfies
\begin{equation}\label{eq:cutcas}
 \mathcal Q_{c,L}=\Pi_L\mathcal Q_c\Pi_L
                                   +2\alpha_LP_L^\partial.
\end{equation}
Even when the infinite discrepancy vanishes,
the finite endpoint terms in these equations remain.
\end{theorem}
\begin{proof}
The truncated raising map sends $e_L$ to zero.
Therefore $[L_L,R_L]=\Pi_L[L,R]\Pi_L-\alpha_LP_L^\partial$,
which proves \eqref{eq:cutdelta}. The purely finite
calculations in Theorem~\ref{thm:brackets} then give
\eqref{eq:cutbracket}. The only path of length two
leaving and returning to the cutoff is
$e_L\to e_{L+1}\to e_L$, with weight $a_L^2$.
Thus
\[
 \Pi_LJ^2\Pi_L=J_L^2+\alpha_LP_L^\partial,\qquad
 \Pi_LB^2\Pi_L=B_L^2+\tfrac14\alpha_LP_L^\partial.
\]
Insert these exact identities in \eqref{eq:W} and
\eqref{eq:Cas} to obtain the two formulas and their signs.
\end{proof}

Every polynomial has an equally explicit correction.
For $0\leq m,n\leq L$ and integer $r\geq0$,
\begin{equation}\label{eq:paths}
 (\Pi_LJ^r\Pi_L-J_L^r)_{mn}
 =\sum_{\substack{v_0=n,\ v_r=m,\ v_j\geq0\\
                  |v_{j+1}-v_j|=1,\ \max_jv_j>L}}
              \prod_{j=0}^{r-1}a_{\min(v_j,v_{j+1})}.
\end{equation}
Matrix multiplication expands into exactly these paths;
the finite power retains precisely those staying at most $L$.
The sum is finite for fixed $r$. For
$P(s_0+iu)=\sum_r\gamma_ru^r$, the correction to
$\Pi_LP(s_0+iJ)\Pi_L$ from $P(s_0+iJ_L)$ is
the same sum multiplied by every exact $\gamma_r$,
where $\gamma=T_d\operatorname{coeff}P$ and $d=\deg P$
(with the zero polynomial handled by zero coefficients).
Thus no polynomial is silently evaluated on a finite
Jacobi matrix in place of compression of the full one.

There is also a useful trace check retaining the whole edge:
\begin{equation}\label{eq:trace}
 \sum_{n=0}^L\delta_n^{(c)}
  =a_L^2-(L+1)(L+2c),\qquad
 \operatorname{tr}\Delta_{c,L}=-(L+1)(L+2c).
\end{equation}
Both follow by telescoping the native recurrence differences.
For $c>0$, the latter is nonzero, so no finite positive
Jacobi cutoff can satisfy all defect-free weight identities.

\section{The actual finite polynomial quotient and its relation defect}

Let $\chi(S)$ be the original full cyclic annihilator of
the sum operator on $E_h^{\otimes k}$, where
$E_h=\C[s]/h$. More explicitly its distinct roots are
$\kappa=\rho_1+\cdots+\rho_k$, with multiplicity
$e_\kappa=\max_{\sum\rho_j=\kappa}(1+\sum_j(m_{\rho_j}-1))$.
This formula follows in the local Chinese-remainder block
$\C[z_1,\ldots,z_k]/(z_1^{m_{\rho_1}},\ldots,z_k^{m_{\rho_k}})$:
the sum operator is $\kappa+\sum_jz_j$,
whose nilpotent part has index exactly
$1+\sum_j(m_{\rho_j}-1)$. Its last nonzero power has
coefficient
$(\sum_j(m_{\rho_j}-1))!/\prod_j(m_{\rho_j}-1)!$
on $\prod_jz_j^{m_{\rho_j}-1}$, which is nonzero.
Taking the maximum over blocks with the same $\kappa$
gives the displayed exponent, with no multiplicity omitted.
Put $q=\deg\chi$, $C_\chi=\C[S]/\chi$, retaining its
original increasing-power remainder coordinates.
The substitution map
$\beta:C_\chi\to E_h^{\otimes k}$,
$[P]\mapsto P(s_1+\cdots+s_k)$, is injective by this
minimal polynomial. The original full physical jet map is
\begin{equation}\label{eq:unitreceiver}
 \eta=\bigl(j_h(g/h)\bigr)^{\otimes k}\beta.
\end{equation}
The multiplier is the complete original local unit in
each factor, including every Taylor coefficient and
root difference. Explicitly, write $H_\rho(s)=h(s)/(s-\rho)^{m_\rho}$
and $u_{\rho,r}=[z^r](g/h)(\rho+z)$ for $0\leq r<m_\rho$.
Then
\begin{equation}\label{eq:unitcoefficients}
 \begin{split}
 u_{\rho,r}
 =\sum_{a=0}^r\frac{g^{(m_\rho+a)}(\rho)}{(m_\rho+a)!}
 \sum_{\substack{\ell_\sigma\geq0\ (\sigma\ne\rho)\\
                 \sum_{\sigma\ne\rho}\ell_\sigma=r-a}}
 \prod_{\sigma\ne\rho}
 \frac{(-1)^{\ell_\sigma}
       \binom{m_\sigma+\ell_\sigma-1}{\ell_\sigma}}
      {(\rho-\sigma)^{m_\sigma+\ell_\sigma}} .
 \end{split}
\end{equation}
This is the product of the Taylor expansion of
$g(\rho+z)/z^{m_\rho}$ with the negative-binomial
expansions of $1/H_\rho(\rho+z)$.
The empty inner product is one when its required sum is
zero and contributes zero otherwise.
In particular
$u_{\rho,0}=g^{(m_\rho)}(\rho)/
 (m_\rho!\prod_{\sigma\ne\rho}(\rho-\sigma)^{m_\sigma})\ne0$,
so multiplication by these complete jets is invertible.
For every original multi-index $0\leq j_a<m_{\rho_a}$ the
full physical map has the explicit entry
\begin{equation}\label{eq:etajets}
 (\eta[P])_{\boldsymbol\rho,\boldsymbol j}
 =\sum_{0\leq\nu_a\leq j_a}
     \left(\prod_{a=1}^k u_{\rho_a,j_a-\nu_a}\right)
          \frac{P^{(|\boldsymbol\nu|)}(\kappa)}
               {\prod_{a=1}^k\nu_a!}.
\end{equation}
The formula follows by expanding
$P(\kappa+z_1+\cdots+z_k)$ with the multinomial theorem.
The full jet of \eqref{eq:physical}
is exactly $\eta[P]$. In particular the algebraic
relation ideal is exactly $\chi\C[S]$, not a
Hilbert orthogonal complement.

\begin{theorem}[Exact failure of grading descent]\label{thm:descent}
For every nonconstant $\chi$, the native polynomial grading
$\mathcal N$ does not preserve $\chi\C[S]$ and hence
does not descend to $C_\chi$. Its first explicit obstruction
is as follows. Expand the original monic polynomial
\begin{equation}\label{eq:chiexpand}
 \chi=p_q+\sum_{j<q}\gamma_jp_j,\qquad
 \gamma_j=\frac{\int_\R\overline{p_j(s_0+iu)}
                \chi(s_0+iu)m(u)\,du}{\omega_j}.
\end{equation}
If at least one $\gamma_j\ne0$, then
\begin{equation}\label{eq:relationdefect1}
 [\mathcal N\chi]_\chi
       =\sum_{j<q}(j-q)\gamma_j[p_j]_\chi\ne0.
\end{equation}
If every $\gamma_j=0$, so $\chi=p_q$, then instead
\begin{equation}\label{eq:relationdefect2}
 [\mathcal N((S-s_0)\chi)]_\chi
       =2\frac{\omega_q}{\omega_{q-1}}[p_{q-1}]_\chi\ne0.
\end{equation}
The displayed relation is present by degree $L=q$ in the
first case and by degree $L=q+1$ in the second.
At $L=q-1$ there are no nonzero polynomial relations.
Consequently $\mathcal W_\lambda$ in
\eqref{eq:Wphysical} descends precisely for $\lambda=0$
among real $\lambda$.
\end{theorem}
\begin{proof}
The $p_j$ for $j<q$ give a basis modulo $\chi$ by their
distinct monic degrees. Applying $\mathcal N-q$ to
\eqref{eq:chiexpand} proves \eqref{eq:relationdefect1};
its nonzero coefficients cannot cancel in that basis.
If $\chi=p_q$, the phase-retaining recurrence
\eqref{eq:physical-polys}--\eqref{eq:recurrence} is
\[
 (S-s_0)p_q=p_{q+1}
                 -\frac{\omega_q}{\omega_{q-1}}p_{q-1}.
\]
Apply $\mathcal N-(q+1)$ to it. Its right side is
$2(\omega_q/\omega_{q-1})p_{q-1}$, which is nonzero
and has degree below $q$. The subtracted multiple
of $(S-s_0)\chi$ is a relation, proving
\eqref{eq:relationdefect2}.
All multiplication and scalar terms of
\eqref{eq:Wphysical} preserve the ideal. The remaining
term is $8\pi\lambda^2\mathcal N$, so for $\lambda\ne0$
the same exact obstruction survives, whereas at zero
the operator is multiplication by $(S-s_0)^2-1/4$.
\end{proof}

For the actual reflection-stable $\chi$ the parity
involution does descend: $\chi(k-S)=(-1)^q\chi(S)$.
Thus $P(S)\mapsto P(k-S)$ preserves the ideal and is
an involution of its quotient. This is compatible
with $p_n(k-S)=(-1)^np_n(S)$.
The even/odd decomposition therefore survives exactly,
although the complete integer grading does not.

\section{Attained compression and both independent boundary terms}

Fix $L\geq q-1$. Write $b_n=[p_n]_\chi$ in the original
remainder coordinates. The native orthonormal reduction
matrix and the original attained form are
\begin{equation}\label{eq:finitequotient}
 \begin{gathered}
 E_L=\left[\frac{i^{-n}b_n}{\sqrt{\omega_n}}\right]_{n=0}^L,
 \qquad K_L^\chi=E_LE_L^*
           =\sum_{n=0}^L\frac{b_nb_n^*}{\omega_n},
 \qquad G_L=(K_L^\chi)^{-1},\\
 R_L^\chi=E_L^*G_L,\qquad
 P_L^\chi=R_L^\chi E_L,\qquad Z_L^\chi=I-P_L^\chi.
 \end{gathered}
\end{equation}
The first $q$ monic columns prove surjectivity.
Hence $R_L^\chi$ is a right inverse, $P_L^\chi$
is the orthogonal projection to $(\ker E_L)^\perp$,
and $\ker E_L$ is precisely the unchanged relation
space expressed in native coordinates.
Also $(R_L^\chi)^*R_L^\chi=G_L$.
To see it is the attained original metric, any lift
of $v$ is $R_L^\chi v+z$ with $E_Lz=0$.
The two terms are orthogonal, and therefore its
squared norm is $v^*G_Lv+\norm z^2$.
This recovers exactly the original full affine minimum.
The phases $i^{-n}$ cancel only in the displayed
diagonal sum for $K_L^\chi$; they remain present
in every off-diagonal operator calculation.

For any polynomial-domain action $A:\PP_L\to\PP_{L+r}$
let $A_{L+r,L}$ be its exact rectangular matrix.
Define its attained finite quotient receiver and relation
defect by
\begin{equation}\label{eq:generalreceiver}
 \widehat A_L=E_{L+r}A_{L+r,L}R_L^\chi,\qquad
 \mathfrak d_{A,L}=E_{L+r}A_{L+r,L}Z_L^\chi.
\end{equation}
Then, for every native coefficient vector $x$,
\begin{equation}\label{eq:relationidentity}
 E_{L+r}A_{L+r,L}x
        =\widehat A_LE_Lx+\mathfrak d_{A,L}x.
\end{equation}
This is the identity $I=P_L^\chi+Z_L^\chi$ with
all maps typed. Thus the non-descent of
Theorem~\ref{thm:descent} has an explicit finite receiver,
not merely a nonidentity.

For a square Hermitian cutoff matrix $A_L$, write
$\overline A_L=E_LA_LR_L^\chi$. Its $G_L$-weighted
matrix is
$G_L\overline A_L=(R_L^\chi)^*A_LR_L^\chi$, so it is
self-adjoint in the attained metric. In particular
$0\leq G_L\overline{\mathsf N}_L\leq L G_L$.
There is no assertion that the eigenvalues of this
compression are integers.
If $A_L,B_L$ are any two square cutoff matrices,
direct multiplication gives the full correction
\begin{equation}\label{eq:compressedbracket}
 [\overline A_L,\overline B_L]
  =E_L[A_L,B_L]R_L^\chi
    -E_LA_LZ_L^\chi B_LR_L^\chi
    +E_LB_LZ_L^\chi A_LR_L^\chi.
\end{equation}
Substitution of \eqref{eq:cutbracket} in this formula
retains the native diagonal discrepancy, the rank-one
cutoff defect, and these two relation terms separately.
They must not be replaced by a rank count.
Writing $E=E_L$, $R=R_L^\chi$, $Z=Z_L^\chi$,
the individual Lie brackets are, explicitly,
\begin{equation}\label{eq:quotientlie}
 \begin{aligned}
 [\overline H_L,\overline E_{\pm,L}]
 ={}&\pm2\overline E_{\pm,L}
   +\tfrac i2E(\Pi_L\Delta_c\Pi_L)R
   -\tfrac i2a_L^2EP_L^\partial R\\
 &-EH_LZE_{\pm,L}R+EE_{\pm,L}ZH_LR,\\
 [\overline E_{+,L},\overline E_{-,L}]
 ={}&\overline H_L/4
       -EE_{+,L}ZE_{-,L}R+EE_{-,L}ZE_{+,L}R .
 \end{aligned}
\end{equation}

The exact compressed prolate matrix obeys the particularly
useful identity
\begin{equation}\label{eq:compressedprolate}
 \begin{split}
 \overline W_{\lambda,L}
 ={}&-\overline J_L^{\,2}
       +8\pi\lambda^2(\overline{\mathsf N}_L+1/4)-1/4\\
 &-E_LJ_LZ_L^\chi J_LR_L^\chi
       -a_L^2E_LP_L^\partial R_L^\chi .
 \end{split}
\end{equation}
To prove it, insert $I=P_L^\chi+Z_L^\chi$
between the two $J_L$ factors of \eqref{eq:cutprolate}
and use $E_LR_L^\chi=I$.
Each of the two subtracted terms is positive
semidefinite in the $G_L$ metric: after multiplying
by $G_L$ they are respectively
$(Z_L^\chi J_LR_L^\chi)^*(Z_L^\chi J_LR_L^\chi)$
and $a_L^2(P_L^\partial R_L^\chi)^*
                           (P_L^\partial R_L^\chi)$.
This fixes both signs and proves an actual attained
quadratic comparison.

For clarity, the uncompressed action of multiplication
by $u=(S-s_0)/i$ does descend, with quotient matrix $M_u$.
The exact rectangular multiplication satisfies
$E_{L+1}J_{L+1,L}=M_uE_L$. Removing its outgoing
edge gives the explicit difference
\begin{equation}\label{eq:Jquotient}
 \overline J_L=M_u-a_L
       \frac{i^{-(L+1)}b_{L+1}}{\sqrt{\omega_{L+1}}}
                                  e_L^*R_L^\chi.
\end{equation}
Equivalently, with every phase multiplied out,
$\overline J_L=M_u+(i/\omega_L)b_{L+1}b_L^*G_L$.
Thus the attained compression of multiplication need
not equal its algebraic quotient action. The outgoing
column, including its phase and mass, is exactly known.

The full rectangular receiver $\widehat W_{\lambda,L}$
from \eqref{eq:generalreceiver} and the square compressed
$\overline W_{\lambda,L}$ are different specified maps.
Their exact bridge is
\begin{equation}\label{eq:fullcompbridge}
 \begin{aligned}
 \widehat W_{\lambda,L}
   &=-M_u^2+8\pi\lambda^2(\overline{\mathsf N}_L+1/4)-1/4,\\
 \overline W_{\lambda,L}
   &=\widehat W_{\lambda,L}
     -\left(\frac{b_{L+1}b_{L-1}^*}{\omega_{L-1}}
             +\frac{b_{L+2}b_L^*}{\omega_L}\right)G_L .
 \end{aligned}
\end{equation}
The first fraction is omitted when $L=0$.
Indeed full multiplication by $u^2$ descends exactly to
$M_u^2$. The only outgoing two-step columns of
$W_\lambda:\PP_L\to\PP_{L+2}$ are
$-a_{L-1}a_Le_{L+1}e_{L-1}^*$
and $-a_La_{L+1}e_{L+2}e_L^*$.
Their $i^{-n}$ phase products change both signs to plus,
and their norm coefficients are $1/\omega_{L-1}$ and
$1/\omega_L$. Subtracting those columns proves the
second formula. This bridge retains outgoing quotient
data as well as the returning-edge correction in
\eqref{eq:compressedprolate}.

All maps can be pushed through the original physical
jet injection \eqref{eq:unitreceiver}: replace $E_L$
by $\eta E_L$ with codomain $\eta(C_\chi)$ and transport
the form by $\eta^{-1}$. This retains the complete
physical unit, not the Euclidean form on a substitute
jet space. For any fixed original observation $\Lambda$
one instead uses the onto map $\Lambda\eta E_L$
to its actual image and repeats the same attained
right-inverse construction there. The identities
\eqref{eq:generalreceiver}--\eqref{eq:compressedbracket}
remain exact, with that observation's actual relation
space; no automatic descent assertion is made for it.

\section{All original phases in explicit compressed matrices}

Set $b_n=[p_n]_\chi$ as above. Direct substitution of
the columns in \eqref{eq:finitequotient} gives
\begin{equation}\label{eq:explicitcompressed}
 \begin{aligned}
 \overline H_L
 &=\left[\sum_{j=0}^{L-1}
       \frac{b_jb_{j+1}^*-b_{j+1}b_j^*}{\omega_j}\right]G_L,\\
 \overline B_L
 &=\left[\frac12\sum_{j=0}^{L-1}
       \frac{b_jb_{j+1}^*+b_{j+1}b_j^*}{\omega_j}\right]G_L,\\
 \overline K_L
 &=\left[\sum_{j=0}^L
                  \frac{(j+c)b_jb_j^*}{\omega_j}\right]G_L,\\
 \overline W_{\lambda,L}
 &=\left[\sum_{j=0}^L\frac{w_jb_jb_j^*}{\omega_j}
      +\sum_{j=0}^{L-2}
             \frac{b_jb_{j+2}^*+b_{j+2}b_j^*}{\omega_j}\right]G_L.
 \end{aligned}
\end{equation}
Every sum with upper endpoint below its lower endpoint
is zero. For example the $j,j+1$ entry of $H$ is
$-ia_j$ and the corresponding phase product is
$i^{-j}i^{j+1}=i$. Their product is $+a_j$,
and $a_j/\sqrt{\omega_j\omega_{j+1}}=1/\omega_j$.
The reversed pair has the negative sign. For $B$
both signs are positive with factor $1/2$.
For the two-step prolate pair the phase product is
$-1$, changing $-a_ja_{j+1}$ to a positive coefficient;
$a_ja_{j+1}/\sqrt{\omega_j\omega_{j+2}}=1/\omega_j$.
These computations prove the formulas while retaining
all actual complex $b_j$ and their cross terms.

\section{Receiver into the existing interpolation-volume determinants}

The original cyclic-volume calculation \cite{CV}%
% BEGIN VERIFIED PUBLIC PROOF LINK
~\href{https://github.com/KokunoYumeto/zeta-function-research-reader/blob/487253b5b01851da5b9b98af4e146fecf1fab902/workbenches/splitzero-tandem/continuations/20260913-toda/tex/cyclic_control_determinant_increments.tex\#L1}{[proof]}%
% END VERIFIED PUBLIC PROOF LINK
{}, CV.1--CV.5,
uses precisely $K_L^\chi$ and $G_L$ from
\eqref{eq:finitequotient}. To prevent collision with the
Hankel determinants $\tau_n$, denote its determinants by
\[
 D_L^\chi=\det K_L^\chi,\qquad
 \widetilde K_L^\chi=K_{L-1}^\chi
                     +b_{L+1}b_{L+1}^*/\omega_{L+1},
 \qquad \widetilde D_L^\chi=\det\widetilde K_L^\chi.
\]
Define $K_{-1}^\chi=0$; a determinant of a singular
matrix is used as such, without an inverse.
The exact existing control identity becomes
\begin{equation}\label{eq:CVreceiver}
 (\epsilon_{h,k,L}^{\rm cyc})^2
 =\underbrace{\frac{\tau_{L+2}(0)\tau_L(0)}
                         {\tau_{L+1}(0)^2}}_{a_L^2}
   \frac{D_{L+1}^\chi+D_{L-1}^\chi
                    -D_L^\chi-\widetilde D_L^\chi}{D_L^\chi}.
\end{equation}
Thus the coefficient of the prolate cutoff edge and
the native Toda norm ratio are exactly the same original
factor multiplying the actual interpolation-volume
increment. The scalar $c$ has not entered that identity.

For a direct verification retaining the complex phase,
put $a=b_L^*G_Lb_L$, $d=b_{L+1}^*G_Lb_{L+1}$ and
$z=b_L^*G_Lb_{L+1}$. The determinant lemma gives
\[
 \frac{D_{L-1}^\chi}{D_L^\chi}=1-\frac a{\omega_L},
 \quad
 \frac{D_{L+1}^\chi}{D_L^\chi}=1+\frac d{\omega_{L+1}},
\]
\[
 \frac{\widetilde D_L^\chi}{D_L^\chi}
 =\left(1-\frac a{\omega_L}\right)
  \left(1+\frac d{\omega_{L+1}}\right)
                +\frac{|z|^2}{\omega_L\omega_{L+1}}.
\]
The rank-two determinant lemma follows by the
Schur determinant of
$\left(\begin{smallmatrix}K&U\\-V^*&I_2\end{smallmatrix}\right)$
in its two orders. Subtracting these displayed ratios
gives $(ad-|z|^2)/(\omega_L\omega_{L+1})$.
Multiplying by $a_L^2=\omega_{L+1}/\omega_L$
proves \eqref{eq:CVreceiver} as the original
$(ad-|z|^2)/\omega_L^2$ control. All formulas hold
at $L=q-1$ as well; then $D_{L-1}^\chi=0$.
For $q=1$ the two-vector Gram determinant is zero.

The same recurrence differences governing the Lie
defect are now exactly
$a_L^2-a_{L-1}^2-2L-2c$, with each of the two
ratios given by \eqref{eq:hankelratios}. Equations
\eqref{eq:todarecurrence}, \eqref{eq:cutprolate},
\eqref{eq:compressedprolate}, and \eqref{eq:CVreceiver}
are their full arithmetic receiving maps.
No zero-location theorem, growing-degree sign, or
asymptotic control follows from the formal Lie
comparison alone.

\begin{thebibliography}{9}\raggedright
\bibitem{CCM} Alain Connes, Caterina Consani and Henri Moscovici,
\emph{Zeta zeros and prolate wave operators: Semilocal adelic operators},
\href{https://arxiv.org/abs/2310.18423v2}{arXiv:2310.18423v2}.
Original author source \texttt{mainc2m24fine.tex}, Sections 2, 3, 6;
Definition \texttt{prolop}, Proposition \texttt{jacobprol},
Theorem \texttt{mainstart}(v), Lemma \texttt{eplusminus1},
and equations \texttt{h}, \texttt{eplusminusricu}, Theorem
\texttt{dzero}. The bracket correction is proved explicitly
above; the uniqueness assertion in \texttt{dzero}(i) is not
imported. Source SHA-256:
\texttt{e0f4b59aef7b4b79366d51fbb81ce18b93bff49aa9e2b63c55b968c470426fcb}.
\bibitem{CV} Split-Zero research programme,
\emph{The cyclic control as an exact change of arithmetic
interpolation volume}, 12 September 2026,
\href{https://github.com/KokunoYumeto/zeta-function-research-reader/blob/487253b5b01851da5b9b98af4e146fecf1fab902/workbenches/splitzero-tandem/continuations/20260913-toda/tex/cyclic_control_determinant_increments.tex\#L1}{\texttt{cyclic\_control\_determinant\_increments.tex}},
CV.1--CV.12, especially CV.1--CV.5 and CV.11--CV.12.
Source SHA-256:
\texttt{06b526ea602e396c00a85bedab0fd59b8170f5e73cc5102c18ab892b86b0359d}.
The formulas used here are independently reproduced
with all source mass and complex cross terms.

% BEGIN VERIFIED PUBLIC PROOF LINK
\href{https://github.com/KokunoYumeto/zeta-function-research-reader/blob/487253b5b01851da5b9b98af4e146fecf1fab902/workbenches/splitzero-tandem/continuations/20260913-toda/tex/cyclic_control_determinant_increments.tex\#L1}{Source proof}.
% END VERIFIED PUBLIC PROOF LINK
\bibitem{Native} \emph{The actual arithmetic multiplier in
Gamma Jacobi coordinates}, 20 September 2026,
\href{https://github.com/KokunoYumeto/zeta-function-research-reader/blob/5992ad50c0f2a06921f1fcaded7b4e38097c0739/workbenches/splitzero-tandem/continuations/20260920-original-kernel-mixed-determinants/providers/NATIVE_JACOBI_COMMUTANT.tex\#L1}{\texttt{NATIVE\_JACOBI\_COMMUTANT.tex}}. The native measure,
original $S$-power coordinate map, and unchanged quotient
are retained; the present native orthogonal basis does
not substitute the Gamma comparison measure.
Source SHA-256:
\texttt{745442cbde7fde3a8177f804816d300f0936f26c921d262af81643c2cc9924c2}.

% BEGIN VERIFIED PUBLIC PROOF LINK
\href{https://github.com/KokunoYumeto/zeta-function-research-reader/blob/5992ad50c0f2a06921f1fcaded7b4e38097c0739/workbenches/splitzero-tandem/continuations/20260920-original-kernel-mixed-determinants/providers/NATIVE_JACOBI_COMMUTANT.tex\#L1}{Source proof}.
% END VERIFIED PUBLIC PROOF LINK
\end{thebibliography}
\end{document}
